
\extrarowsep=.8mm
\begin{tabu}to .95\linewidth{X[c]X[c,1.2]X[c,.6]X[c,1.15]}
	\multicolumn{4}{c}{\zihao{-2}\bf 西南大学本科毕业论文（设计）开题报告}\\\tabucline-
	\multicolumn{1}{|c|}{\bf 论文(设计)题目} & \multicolumn{3}{c|}{\fangsong 本科毕业论文模板---附带简单的说明}  \\\hline
		\multicolumn{1}{|c|}{\bf\ziju{1.2}{学生姓名}} & \multicolumn{1}{c|}{\fangsong 包小敏} &  \multicolumn{1}{c|}{\bf\ziju{2}{学号}} & \multicolumn{1}{c|}{\fangsong 22222222} \\\hline
	\multicolumn{1}{|c|}{\bf 学院\,/\,专业\,/\,年级} & \multicolumn{3}{c|}{\fangsong 数学与统计学院/数学与应用数学专业/2021级}  \\\hline
	\multicolumn{1}{|c|}{\bf\ziju{1.2}{指导教师}} & \multicolumn{1}{c|}{\fangsong 奈斯如拉.阿布都艾山} &  \multicolumn{1}{c|}{\bf 评阅时间} & \multicolumn{1}{c|}{\today} \\\tabucline-
\tabuphantomline
\end{tabu}
\vspace{-0.9mm}

\begin{tabu}to 0.95\linewidth{|X|}
	\parbox[t][9.8cm][t]{14.5cm}{
		1. 本课题研究意义：\fangsong%
		%具体评阅意见也可在此引入
		具体内容在文件夹{\tt table} 中的文件 kaiti 1.tex 内填写。
		\zcomments
	}\\\tabucline-		
	\tabuphantomline
	\parbox[t][9.1cm][t]{14.5cm}{
		2. 研究内容：\fangsong 具体内容在文件夹~{\tt table} 中的文件~{\tt kaiti\_1.tex} 内这个位置填写。
		%\zcomments
		}
	\\\tabucline-		
	%\tabuphantomline
\end{tabu}	
